\section{讨论与结论}

我们提出了 \ours{}，一个沙箱化评估平台与基准，覆盖八个企业生产力领域中的 1,150 个专家策划任务，包含 512 个工具和由主题专家编写的 SQL 验证器。实验揭示了若干我们认为对企业级 LLM 智能体发展具有广泛意义的发现。

\textbf{当前智能体远未达到企业就绪。}
即使在预言式工具访问这一最有利的检索设置下，最佳模型的任务成功率也仅为 37.4\%；性能还会随跨度长度和跨领域耦合单调下降。关键在于，失败模式并非随机：我们的定性和定量分析表明，智能体在\textit{权限与流程合规}方面最为困难，例如遵守策略和处理级联状态转移，而不是基本任务完成。此外，即便前沿模型也仅约一半时间能可靠拒绝不可行任务，最佳值为 53.9\%，远低于无人监督部署所需的鲁棒性。这些结果说明，仅靠扩展模型容量无法弥合企业可靠性差距。

\textbf{规划是主要瓶颈，而非工具执行。}
我们的计划条件消融表明，人工编写计划能在跨模型、跨领域时带来 14--35 个百分点的提升，远大于自动规划或更复杂多智能体编排的收益。值得注意的是，以人工计划为条件的 Qwen3-4B 可与同条件下大得多的模型竞争，说明一旦战略性推理被外化，即使小模型也能忠实执行。这种规划与执行能力的分离意味着，核心挑战是约束感知的计划生成，而非工具调用熟练度。与此一致，干扰工具不会显著进一步伤害性能，这再次确认工具检索不是约束性能的主要因素。因此，推进长程、策略感知的规划，是提升 \ours{} 上智能体性能的最高杠杆方向。

\textbf{思考预算很重要，但存在领域特异的上限。}
增加测试时计算会在多数领域显著提升性能，但扩展并非普遍单调：一些领域较早进入平台期，说明额外推理 token 不能弥补领域知识或策略理解的根本缺口。未来工作应研究如何更自适应地分配测试时计算，以及针对高约束领域的定向训练能否提高这些上限。

\textbf{未来方向。}
我们的结果提出三项具体研究优先级。第一，\textit{约束感知的计划生成}：在提交动作序列之前，显式推理策略约束、副作用依赖和先决条件结构的方法。第二，\textit{长程状态管理}：在多次工具调用中维护连贯世界状态的机制，例如情景记忆或结构化状态表示，以避免我们观察到的、随跨度增长而累积的错误。第三，\textit{安全拒绝与升级}：智能体必须能可靠识别不可行或违反策略的请求并干净地弃权，而这一能力在当前所有受评系统中仍然较弱。

我们将发布 \ours{}、其沙箱环境和评估工具，以支持开放、社区驱动的研究。沙箱采用模块化、可扩展设计，可随企业实践演变加入新的领域、工具和工作流。通过将智能体评测扎根于真实、约束密集的企业工作流，\ours{} 旨在促使领域关注那些真正决定 LLM 智能体能否作为可靠\textit{AI 员工}部署的规划、安全和策略合规能力。
